\documentclass[12pt,a4paper]{article}
\usepackage[utf8]{vietnam}
\usepackage{amsmath,amssymb}
\usepackage{tabularx}
\usepackage{tikz}
\usepackage{xcolor}
\usepackage[top=0.6cm, bottom=0.6cm, left=0.6cm, right=0.6cm]{geometry}
\usepackage[dethi]{ex_test}

\setlength{\parskip}{0pt}
\setlength{\parindent}{0pt}

% Dòng kẻ chấm mờ đều đặn chuẩn viết tay (khoảng cách 20.5pt)
\newcommand{\dongke}[1]{%
  \begingroup%
  \setlength{\parskip}{0pt}%
  \baselineskip=20.5pt%
  \vspace{1mm}%
  \foreach \i in {1,...,#1}{%
    \noindent{\color{gray!60}\dotfill}\par%
  }%
  \vspace{0.5mm}%
  \endgroup%
}

% Định nghĩa ô bài tập tự luận trong bảng (dành cho Câu 1 Trang 1)
\newcommand{\oBaiTap}[3][6]{%
  \vspace{1.2mm}%
  \noindent\textbf{#2} #3\par%
  \dongke{#1}%
}

% Câu trắc nghiệm 1 dòng: 4 đáp án dàn hàng ngang
\newcommand{\cauTN}[6]{% \cauTN{số câu}{đề}{A}{B}{C}{D}
  \noindent\textbf{Câu #1.} #2\par\vspace{0.3mm}%
  \noindent%
  \makebox[0.25\linewidth][l]{\textbf{A.} #3}%
  \makebox[0.25\linewidth][l]{\textbf{B.} #4}%
  \makebox[0.25\linewidth][l]{\textbf{C.} #5}%
  \makebox[0.25\linewidth][l]{\textbf{D.} #6}\par\vspace{0.6mm}%
}

% Câu trắc nghiệm 2 dòng: mỗi dòng 2 đáp án (dành cho câu có đáp án dài)
\newcommand{\cauTNB}[6]{% \cauTNB{số câu}{đề}{A}{B}{C}{D}
  \noindent\textbf{Câu #1.} #2\par\vspace{0.3mm}%
  \noindent%
  \makebox[0.5\linewidth][l]{\textbf{A.} #3}%
  \makebox[0.5\linewidth][l]{\textbf{B.} #4}\par\vspace{0.3mm}%
  \noindent%
  \makebox[0.5\linewidth][l]{\textbf{C.} #5}%
  \makebox[0.5\linewidth][l]{\textbf{D.} #6}\par\vspace{0.6mm}%
}

\begin{document}
\pagestyle{empty}
% =============================================================================
% ========================= TOÁN 9 - TUẦN 05 - BUỔI 02: MÃ ĐỀ 101 =============
% ========================= MẶT TRƯỚC (TRANG 1/2) ==============================
% =============================================================================

\noindent
\begin{tabularx}{\linewidth}{|X|c|}
\hline
\textbf{PHIẾU KIỂM TRA ĐỊNH KỲ TOÁN 9 - TUẦN 05 - BUỔI 02} & \textbf{Mã đề: 101 - Trang 1/2} \\
\textit{Chủ đề: Hệ phương trình - Giải tam giác vuông - Lập hệ thực tế} & \textit{Thời gian: 45 phút} \\
\hline
\end{tabularx}

\vspace{1mm}

\noindent
\begin{tabularx}{\linewidth}{|X|}
\hline
\rule[-3.5mm]{0pt}{10mm}\textbf{Họ và tên học sinh:} \hfill \textbf{Mã HS:} \makebox[3.5cm]{\dotfill} \\
\hline
\end{tabularx}

\vspace{1mm}

\noindent
\begin{tabularx}{\linewidth}{|p{3.2cm}|X|}
\hline
\centering\textbf{ĐIỂM SỐ} & \centering\textbf{LỜI NHẬN XÉT CỦA THẦY CÔ} \tabularnewline
\hline
\rule{0pt}{1.8cm} & \tabularnewline
\hline
\end{tabularx}

\vspace{1.5mm}

\noindent\textbf{\large PHẦN I. TRẮC NGHIỆM KHÁCH QUAN (2,5 điểm)} \hfill \textit{(Khoanh tròn đáp án đúng)}\par
\vspace{0.8mm}
\cauTN{1}{Cho tam giác $ABC$ vuông tại $A$. Khẳng định nào sau đây đúng?}{$AC = BC \cdot \sin B$}{$AC = BC \cdot \cos B$}{$AC = AB \cdot \sin B$}{$AC = BC \cdot \tan B$}
\cauTN{2}{Cho tam giác $ABC$ vuông tại $A$, có $BC = 10\text{ cm}$ và $\widehat{B} = 30^\circ$. Độ dài cạnh $AC$ bằng:}{$5\sqrt{3}\text{ cm}$}{$5\text{ cm}$}{$10\sqrt{3}\text{ cm}$}{$20\text{ cm}$}
\cauTN{3}{Cho tam giác $ABC$ vuông tại $A$, có $AB = 6\text{ cm}$ và $\cos B = 0{,}6$. Độ dài cạnh huyền $BC$ bằng:}{$3{,}6\text{ cm}$}{$10\text{ cm}$}{$8\text{ cm}$}{$12\text{ cm}$}
\cauTN{4}{Cho tam giác $ABC$ vuông tại $A$, có $AB = 5\text{ cm}$ và $\widehat{B} = 45^\circ$. Độ dài cạnh $AC$ bằng:}{$5\sqrt{2}\text{ cm}$}{$5\text{ cm}$}{$10\text{ cm}$}{$2{,}5\text{ cm}$}
\cauTN{5}{Một chiếc thang dài $4\text{ m}$ tựa vào tường, tạo với mặt đất góc $60^\circ$. Khoảng cách từ chân thang đến tường là:}{$2\text{ m}$}{$2\sqrt{3}\text{ m}$}{$4\sqrt{3}\text{ m}$}{$8\text{ m}$}
\vspace{1mm}

\noindent\textbf{\large PHẦN II. TỰ LUẬN (7,5 điểm)}

\noindent\textbf{Câu 1 (2,5 điểm). Giải các hệ phương trình sau:}
\vspace{0.6mm}

\noindent
\begin{tabularx}{\linewidth}{|X|X|}
\hline
\oBaiTap[16]{a)}{$\begin{cases} 3x - 2y = 7 \\ 2x + 3y = 9 \end{cases}$} & 
\oBaiTap[16]{b)}{$\begin{cases} 2x + 3y = 7 \\ 3x - 2y = 4 \end{cases}$} \\
\hline
\end{tabularx}

\newpage
% =============================================================================
% ========================= TOÁN 9 - TUẦN 05 - BUỔI 02: MÃ ĐỀ 101 =============
% ========================= MẶT SAU (TRANG 2/2) ================================
% =============================================================================

\noindent
\begin{tabularx}{\linewidth}{|X|c|}
\hline
\textbf{PHẦN II. TỰ LUẬN (Tiếp theo) - TOÁN 9 TUẦN 05 BUỔI 02} & \textbf{Mã đề: 101 - Trang 2/2} \\
\hline
\end{tabularx}

\vspace{1.5mm}

\noindent\textbf{\large Câu 2 (2,5 điểm). Giải tam giác vuông:} Cho tam giác $ABC$ vuông tại $A$, biết cạnh $AB = 6\text{ cm}$ và $\cos B = \dfrac{3}{5}$.\par
\vspace{0.4mm}
$a)$ Tính độ dài cạnh huyền $BC$ và cạnh góc vuông $AC$.\par
\vspace{0.3mm}
$b)$ Kẻ đường cao $AH$ vuông góc với $BC$ ($H \in BC$). Tính độ dài đoạn thẳng $AH$ và đoạn thẳng $BH$.\par
\vspace{0.8mm}

\noindent
\begin{tabularx}{\linewidth}{|X|X|}
\hline
\dongke{13} & \dongke{13} \\
\hline
\end{tabularx}

\vspace{1.5mm}

\noindent\textbf{\large Câu 3 (2,5 điểm). Bài toán thực tế lập hệ phương trình:} Nhân dịp khai giảng năm học mới, bạn Nam đến nhà sách mua một chiếc cặp sách và một bộ dụng cụ học tập với tổng giá niêm yết là $300\,000$ đồng. Do có chương trình khuyến mãi, nhà sách giảm giá $10\%$ cho chiếc cặp sách và giảm giá $20\%$ cho bộ dụng cụ học tập, vì vậy bạn Nam chỉ phải trả tổng cộng $258\,000$ đồng. Tính giá tiền niêm yết ban đầu của chiếc cặp sách và bộ dụng cụ học tập.\par
\vspace{0.8mm}

\noindent
\begin{tabularx}{\linewidth}{|X|X|}
\hline
\dongke{16} & \dongke{16} \\
\hline
\end{tabularx}

\newpage
% =============================================================================
% ========================= TOÁN 9 - TUẦN 05 - BUỔI 02: MÃ ĐỀ 102 =============
% ========================= MẶT TRƯỚC (TRANG 1/2) ==============================
% =============================================================================

\noindent
\begin{tabularx}{\linewidth}{|X|c|}
\hline
\textbf{PHIẾU KIỂM TRA ĐỊNH KỲ TOÁN 9 - TUẦN 05 - BUỔI 02} & \textbf{Mã đề: 102 - Trang 1/2} \\
\textit{Chủ đề: Hệ phương trình - Giải tam giác vuông - Lập hệ thực tế} & \textit{Thời gian: 45 phút} \\
\hline
\end{tabularx}

\vspace{1mm}

\noindent
\begin{tabularx}{\linewidth}{|X|}
\hline
\rule[-3.5mm]{0pt}{10mm}\textbf{Họ và tên học sinh:} \hfill \textbf{Mã HS:} \makebox[3.5cm]{\dotfill} \\
\hline
\end{tabularx}

\vspace{1mm}

\noindent
\begin{tabularx}{\linewidth}{|p{3.2cm}|X|}
\hline
\centering\textbf{ĐIỂM SỐ} & \centering\textbf{LỜI NHẬN XÉT CỦA THẦY CÔ} \tabularnewline
\hline
\rule{0pt}{1.8cm} & \tabularnewline
\hline
\end{tabularx}

\vspace{1.5mm}

\noindent\textbf{\large PHẦN I. TRẮC NGHIỆM KHÁCH QUAN (2,5 điểm)} \hfill \textit{(Khoanh tròn đáp án đúng)}\par
\vspace{0.8mm}
\cauTN{1}{Cho tam giác $ABC$ vuông tại $A$. Khẳng định nào sau đây đúng?}{$AB = BC \cdot \sin B$}{$AB = BC \cdot \cos B$}{$AB = AC \cdot \cos B$}{$AB = BC \cdot \cot B$}
\cauTN{2}{Cho tam giác $ABC$ vuông tại $A$, có $BC = 12\text{ cm}$ và $\widehat{C} = 60^\circ$. Độ dài cạnh $AC$ bằng:}{$6\text{ cm}$}{$6\sqrt{3}\text{ cm}$}{$4\sqrt{3}\text{ cm}$}{$24\text{ cm}$}
\cauTN{3}{Cho tam giác $ABC$ vuông tại $A$, có $AC = 8\text{ cm}$ và $\sin B = 0{,}8$. Độ dài cạnh huyền $BC$ bằng:}{$10\text{ cm}$}{$6{,}4\text{ cm}$}{$6\text{ cm}$}{$16\text{ cm}$}
\cauTN{4}{Cho tam giác $ABC$ vuông tại $A$, có $AB = 4\text{ cm}$ và $\tan B = 1{,}5$. Độ dài cạnh $AC$ bằng:}{$6\text{ cm}$}{$\dfrac{8}{3}\text{ cm}$}{$5\text{ cm}$}{$8\text{ cm}$}
\cauTN{5}{Một chiếc thang dài $5\text{ m}$ tựa vào tường, chân thang cách tường $2{,}5\text{ m}$. Góc tạo bởi thang và mặt đất là:}{$30^\circ$}{$60^\circ$}{$45^\circ$}{$75^\circ$}
\vspace{1mm}

\noindent\textbf{\large PHẦN II. TỰ LUẬN (7,5 điểm)}

\noindent\textbf{Câu 1 (2,5 điểm). Giải các hệ phương trình sau:}
\vspace{0.6mm}

\noindent
\begin{tabularx}{\linewidth}{|X|X|}
\hline
\oBaiTap[16]{a)}{$\begin{cases} 2x - 3y = 1 \\ 3x + 2y = 8 \end{cases}$} & 
\oBaiTap[16]{b)}{$\begin{cases} 3x + 2y = 11 \\ 2x - 3y = 3 \end{cases}$} \\
\hline
\end{tabularx}

\newpage
% =============================================================================
% ========================= TOÁN 9 - TUẦN 05 - BUỔI 02: MÃ ĐỀ 102 =============
% ========================= MẶT SAU (TRANG 2/2) ================================
% =============================================================================

\noindent
\begin{tabularx}{\linewidth}{|X|c|}
\hline
\textbf{PHẦN II. TỰ LUẬN (Tiếp theo) - TOÁN 9 TUẦN 05 BUỔI 02} & \textbf{Mã đề: 102 - Trang 2/2} \\
\hline
\end{tabularx}

\vspace{1.5mm}

\noindent\textbf{\large Câu 2 (2,5 điểm). Giải tam giác vuông:} Cho tam giác $ABC$ vuông tại $A$, biết cạnh $AC = 12\text{ cm}$ và $\sin B = \dfrac{3}{5}$.\par
\vspace{0.4mm}
$a)$ Tính độ dài cạnh huyền $BC$ và cạnh góc vuông $AB$.\par
\vspace{0.3mm}
$b)$ Kẻ đường cao $AH$ vuông góc với $BC$ ($H \in BC$). Tính độ dài đoạn thẳng $AH$ và diện tích tam giác $ABC$.\par
\vspace{0.8mm}

\noindent
\begin{tabularx}{\linewidth}{|X|X|}
\hline
\dongke{13} & \dongke{13} \\
\hline
\end{tabularx}

\vspace{1.5mm}

\noindent\textbf{\large Câu 3 (2,5 điểm). Bài toán thực tế lập hệ phương trình:} Nhân dịp năm học mới, bạn Bình mua một chiếc áo đồng phục và một đôi giày thể thao với tổng giá niêm yết là $500\,000$ đồng. Cửa hàng giảm giá $15\%$ cho chiếc áo đồng phục và giảm giá $10\%$ cho đôi giày thể thao, do đó bạn Bình chỉ phải trả tổng cộng $435\,000$ đồng. Tính giá tiền niêm yết ban đầu của mỗi sản phẩm.\par
\vspace{0.8mm}

\noindent
\begin{tabularx}{\linewidth}{|X|X|}
\hline
\dongke{16} & \dongke{16} \\
\hline
\end{tabularx}

\newpage
% =============================================================================
% ========================= TOÁN 9 - TUẦN 05 - BUỔI 02: MÃ ĐỀ 103 =============
% ========================= MẶT TRƯỚC (TRANG 1/2) ==============================
% =============================================================================

\noindent
\begin{tabularx}{\linewidth}{|X|c|}
\hline
\textbf{PHIẾU KIỂM TRA ĐỊNH KỲ TOÁN 9 - TUẦN 05 - BUỔI 02} & \textbf{Mã đề: 103 - Trang 1/2} \\
\textit{Chủ đề: Hệ phương trình - Giải tam giác vuông - Lập hệ thực tế} & \textit{Thời gian: 45 phút} \\
\hline
\end{tabularx}

\vspace{1mm}

\noindent
\begin{tabularx}{\linewidth}{|X|}
\hline
\rule[-3.5mm]{0pt}{10mm}\textbf{Họ và tên học sinh:} \hfill \textbf{Mã HS:} \makebox[3.5cm]{\dotfill} \\
\hline
\end{tabularx}

\vspace{1mm}

\noindent
\begin{tabularx}{\linewidth}{|p{3.2cm}|X|}
\hline
\centering\textbf{ĐIỂM SỐ} & \centering\textbf{LỜI NHẬN XÉT CỦA THẦY CÔ} \tabularnewline
\hline
\rule{0pt}{1.8cm} & \tabularnewline
\hline
\end{tabularx}

\vspace{1.5mm}

\noindent\textbf{\large PHẦN I. TRẮC NGHIỆM KHÁCH QUAN (2,5 điểm)} \hfill \textit{(Khoanh tròn đáp án đúng)}\par
\vspace{0.8mm}
\cauTN{1}{Cho tam giác $ABC$ vuông tại $A$. Khẳng định nào sau đây SAI?}{$AC = BC \cdot \sin B$}{$AB = BC \cdot \cos B$}{$AC = AB \cdot \cot B$}{$AC = AB \cdot \tan B$}
\cauTN{2}{Cho tam giác $ABC$ vuông tại $A$, có $BC = 8\text{ cm}$ và $\widehat{B} = 60^\circ$. Độ dài cạnh $AB$ bằng:}{$4\sqrt{3}\text{ cm}$}{$4\text{ cm}$}{$8\sqrt{3}\text{ cm}$}{$16\text{ cm}$}
\cauTN{3}{Cho tam giác $ABC$ vuông tại $A$, có $AB = 9\text{ cm}$ và $\cos B = \dfrac{3}{5}$. Độ dài cạnh huyền $BC$ bằng:}{$5{,}4\text{ cm}$}{$12\text{ cm}$}{$15\text{ cm}$}{$25\text{ cm}$}
\cauTN{4}{Cho tam giác $ABC$ vuông tại $A$, có $AB = 6\text{ cm}$ và $\widehat{B} = 45^\circ$. Độ dài cạnh $AC$ bằng:}{$6\sqrt{2}\text{ cm}$}{$6\text{ cm}$}{$3\text{ cm}$}{$12\text{ cm}$}
\cauTN{5}{Một cột cờ cao $7\text{ m}$ có bóng trên mặt đất dài $7\text{ m}$. Góc tạo bởi tia nắng và mặt đất là:}{$45^\circ$}{$30^\circ$}{$60^\circ$}{$50^\circ$}
\vspace{1mm}

\noindent\textbf{\large PHẦN II. TỰ LUẬN (7,5 điểm)}

\noindent\textbf{Câu 1 (2,5 điểm). Giải các hệ phương trình sau:}
\vspace{0.6mm}

\noindent
\begin{tabularx}{\linewidth}{|X|X|}
\hline
\oBaiTap[16]{a)}{$\begin{cases} 3x + 4y = 10 \\ 2x - 3y = 1 \end{cases}$} & 
\oBaiTap[16]{b)}{$\begin{cases} 2x + 3y = 5 \\ 3x - 2y = 1 \end{cases}$} \\
\hline
\end{tabularx}

\newpage
% =============================================================================
% ========================= TOÁN 9 - TUẦN 05 - BUỔI 02: MÃ ĐỀ 103 =============
% ========================= MẶT SAU (TRANG 2/2) ================================
% =============================================================================

\noindent
\begin{tabularx}{\linewidth}{|X|c|}
\hline
\textbf{PHẦN II. TỰ LUẬN (Tiếp theo) - TOÁN 9 TUẦN 05 BUỔI 02} & \textbf{Mã đề: 103 - Trang 2/2} \\
\hline
\end{tabularx}

\vspace{1.5mm}

\noindent\textbf{\large Câu 2 (2,5 điểm). Giải tam giác vuông:} Cho tam giác $ABC$ vuông tại $A$, biết cạnh $AB = 5\text{ cm}$ và $\tan B = \dfrac{12}{5}$.\par
\vspace{0.4mm}
$a)$ Tính độ dài cạnh góc vuông $AC$ và cạnh huyền $BC$.\par
\vspace{0.3mm}
$b)$ Kẻ đường cao $AH$ vuông góc với $BC$ ($H \in BC$). Tính độ dài đoạn thẳng $AH$ và tỉ số lượng giác $\sin C$.\par
\vspace{0.8mm}

\noindent
\begin{tabularx}{\linewidth}{|X|X|}
\hline
\dongke{13} & \dongke{13} \\
\hline
\end{tabularx}

\vspace{1.5mm}

\noindent\textbf{\large Câu 3 (2,5 điểm). Bài toán thực tế lập hệ phương trình:} Để chuẩn bị góc học tập, bạn Hoa mua một chiếc máy tính cầm tay và một chiếc đèn bàn học với tổng giá niêm yết là $600\,000$ đồng. Nhân dịp khuyến mãi, cửa hàng giảm giá $20\%$ cho máy tính cầm tay và giảm giá $10\%$ cho chiếc đèn bàn, nên bạn Hoa chỉ phải trả tổng cộng $510\,000$ đồng. Tính giá tiền niêm yết ban đầu của mỗi mặt hàng.\par
\vspace{0.8mm}

\noindent
\begin{tabularx}{\linewidth}{|X|X|}
\hline
\dongke{16} & \dongke{16} \\
\hline
\end{tabularx}

\newpage
% =============================================================================
% ========================= TOÁN 9 - TUẦN 05 - BUỔI 02: MÃ ĐỀ 104 =============
% ========================= MẶT TRƯỚC (TRANG 1/2) ==============================
% =============================================================================

\noindent
\begin{tabularx}{\linewidth}{|X|c|}
\hline
\textbf{PHIẾU KIỂM TRA ĐỊNH KỲ TOÁN 9 - TUẦN 05 - BUỔI 02} & \textbf{Mã đề: 104 - Trang 1/2} \\
\textit{Chủ đề: Hệ phương trình - Giải tam giác vuông - Lập hệ thực tế} & \textit{Thời gian: 45 phút} \\
\hline
\end{tabularx}

\vspace{1mm}

\noindent
\begin{tabularx}{\linewidth}{|X|}
\hline
\rule[-3.5mm]{0pt}{10mm}\textbf{Họ và tên học sinh:} \hfill \textbf{Mã HS:} \makebox[3.5cm]{\dotfill} \\
\hline
\end{tabularx}

\vspace{1mm}

\noindent
\begin{tabularx}{\linewidth}{|p{3.2cm}|X|}
\hline
\centering\textbf{ĐIỂM SỐ} & \centering\textbf{LỜI NHẬN XÉT CỦA THẦY CÔ} \tabularnewline
\hline
\rule{0pt}{1.8cm} & \tabularnewline
\hline
\end{tabularx}

\vspace{1.5mm}

\noindent\textbf{\large PHẦN I. TRẮC NGHIỆM KHÁCH QUAN (2,5 điểm)} \hfill \textit{(Khoanh tròn đáp án đúng)}\par
\vspace{0.8mm}
\cauTN{1}{Cho tam giác $ABC$ vuông tại $A$. Khẳng định nào sau đây đúng?}{$AC = AB \cdot \cot B$}{$AB = AC \cdot \tan B$}{$AC = AB \cdot \tan B$}{$AC = BC \cdot \cos B$}
\cauTN{2}{Cho tam giác $ABC$ vuông tại $A$, có $BC = 14\text{ cm}$ và $\widehat{B} = 30^\circ$. Độ dài cạnh $AC$ bằng:}{$7\sqrt{3}\text{ cm}$}{$14\sqrt{3}\text{ cm}$}{$7\text{ cm}$}{$28\text{ cm}$}
\cauTN{3}{Cho tam giác $ABC$ vuông tại $A$, có $AC = 12\text{ cm}$ và $\sin B = \dfrac{3}{5}$. Độ dài cạnh huyền $BC$ bằng:}{$7{,}2\text{ cm}$}{$20\text{ cm}$}{$16\text{ cm}$}{$25\text{ cm}$}
\cauTN{4}{Cho tam giác $ABC$ vuông tại $A$, có $AB = 8\text{ cm}$ và $\tan B = 0{,}75$. Độ dài cạnh $AC$ bằng:}{$10\text{ cm}$}{$8\text{ cm}$}{$6\text{ cm}$}{$4\text{ cm}$}
\cauTN{5}{Một chiếc thang dài $6\text{ m}$ tựa vào tường, tạo với mặt đất góc $60^\circ$. Khoảng cách từ chân thang đến tường là:}{$3\sqrt{3}\text{ m}$}{$6\sqrt{3}\text{ m}$}{$3\text{ m}$}{$12\text{ m}$}
\vspace{1mm}

\noindent\textbf{\large PHẦN II. TỰ LUẬN (7,5 điểm)}

\noindent\textbf{Câu 1 (2,5 điểm). Giải các hệ phương trình sau:}
\vspace{0.6mm}

\noindent
\begin{tabularx}{\linewidth}{|X|X|}
\hline
\oBaiTap[16]{a)}{$\begin{cases} 4x - 3y = 5 \\ 3x + 2y = 8 \end{cases}$} & 
\oBaiTap[16]{b)}{$\begin{cases} 3x - 2y = 2 \\ 2x + 3y = 10 \end{cases}$} \\
\hline
\end{tabularx}

\newpage
% =============================================================================
% ========================= TOÁN 9 - TUẦN 05 - BUỔI 02: MÃ ĐỀ 104 =============
% ========================= MẶT SAU (TRANG 2/2) ================================
% =============================================================================

\noindent
\begin{tabularx}{\linewidth}{|X|c|}
\hline
\textbf{PHẦN II. TỰ LUẬN (Tiếp theo) - TOÁN 9 TUẦN 05 BUỔI 02} & \textbf{Mã đề: 104 - Trang 2/2} \\
\hline
\end{tabularx}

\vspace{1.5mm}

\noindent\textbf{\large Câu 2 (2,5 điểm). Giải tam giác vuông:} Cho tam giác $ABC$ vuông tại $A$, biết cạnh huyền $BC = 15\text{ cm}$ và $\cos B = \dfrac{4}{5}$.\par
\vspace{0.4mm}
$a)$ Tính độ dài hai cạnh góc vuông $AB$ và $AC$.\par
\vspace{0.3mm}
$b)$ Kẻ đường cao $AH$ vuông góc với $BC$ ($H \in BC$). Tính độ dài đoạn thẳng $AH$ và diện tích tam giác $ABC$.\par
\vspace{0.8mm}

\noindent
\begin{tabularx}{\linewidth}{|X|X|}
\hline
\dongke{13} & \dongke{13} \\
\hline
\end{tabularx}

\vspace{1.5mm}

\noindent\textbf{\large Câu 3 (2,5 điểm). Bài toán thực tế lập hệ phương trình:} Bạn Minh mua một chiếc ba lô và một chiếc mũ bảo hiểm với tổng giá niêm yết là $400\,000$ đồng. Cửa hàng giảm giá $10\%$ cho chiếc ba lô và giảm giá $15\%$ cho chiếc mũ bảo hiểm, nên bạn Minh chỉ phải trả tổng cộng $352\,000$ đồng. Tính giá tiền niêm yết ban đầu của mỗi mặt hàng.\par
\vspace{0.8mm}

\noindent
\begin{tabularx}{\linewidth}{|X|X|}
\hline
\dongke{16} & \dongke{16} \\
\hline
\end{tabularx}

\end{document}
